\chapter{Test 6: Congestion control and Reliability}
In this section we study the effect of heterogeneous bandwidth with different
packet loss rates in different configurations our protocol implementation supports.

\section{Evaluation scenario}
To evaluate the performance under the assumptions given in this task four clients
were assigned a IP address from loopback address block (127.0.0.0/8). The server
belongs also to this address block. This network setup of ours is given in
figure \ref{fig:task5-setup}.  These additional addresses were created as
alises with command:\\
\\
\texttt{sudo ifconfig lo0 alias 127.0.0.2}
\\

Dummynet was used to create link characteristics between client and server
machines with e.g. following command for each link between client and server
(naturally changing the IP addresses as required): \\ 
\\
\texttt{sudo ipfw add pipe 1 ip from 127.0.0.2 to 127.0.0.10; \\sudo ipfw pipe 1 config bw 500kb }
\\
The packet loss was created with our transport layer implementation which is
capable to simulate artificial packet loss and delays. It's good to notice
however that this is not the most realistic situation as packet loss is only
happening on the link from server to client. This means that when a weakly
realiabile subscriptions are used all acknowledgements are not experiencing any
packet loss except in case that happens naturally. The loss rates for clients
are 0.15 and 0.05. The losses are bursty meaning that when a loss happens it's
more likely that also the next message will be lost.

In comparison to high traffic volume sensors the results for temperature sensors
are given in appendix \ref{appendix:task6}.

\section{Results}
The scenario where we have no weak-reliability guarantees or active congestion
control are shown in table \ref{table:task6-camera-a} for camera sensor data. We
observe quite high packet loss rate due to fragmentation while also the effect
of bandwidth is clearly shown.

The results when weakly reliable connections are used are shown in table
\ref{table:task6-camera-b}. We see no improvement in packet losses which is
likely due to high traffic volume of the camera sensors. Because of this the
sensor updates which are lost are going to be retransmitted while due to burstly
updates from camera most likely cause the messages to be short lived.
Retransmitting short lived messages like camera updates are bad thing with our
current implementation because this blocks the transmission of newer updates as
the server is not threaded.

The statistics when message skipping is turned on are shown in table
\ref{table:task6-camera-c}. As assumed the skipping does not provide any good
mean to prevent the congestion effects. The statistics when message skipping is turned on with weakly reliable transport are shown in table
\ref{table:task6-camera-d}.

The results for temperature sensors (shown in appendix \ref{appendix:task6})
show that the reliable connections are working properly (packet loss is
minimized) when the communication flow has low traffic intensity.


\begin{table}
\centering
\begin{tabular}{|c|c|c|c|c|c|c|c|c|}
\hline
client & E[D] & max(D) & min(D) & Var(D) & E[s] & max(s) & Var(s) & plr\\
\hline
500kbs 0.05 & 0.044 & 0.080 & 0.000 & 0.000 & 0.390 & 2.000 & 0.345 & 0.16\\
\hline
500kbs 0.15 & 0.044 & 0.070 & 0.000 & 0.000 & 2.203 & 13.000 &9.751 & 0 0.5 \\
\hline
100kbs 0.05 & 2.336 & 4.350 & 0.100 & 1.523 & 8.407 & 25.000 & 65.970 & 0.35 \\
\hline
100kbs 0.15 & 1.323 & 2.590 & 0.020 & 0.506 & 5.085 & 21.000 & 34.838 & 0.55 \\
\hline
\end{tabular}
\caption{Task 6: Receive statistics for camera sensor data (no realibility guarantees, no congestion control)}
\label{table:task6-camera-a}
\end{table}

\begin{table}
\centering
\begin{tabular}{|c|c|c|c|c|c|c|c|c|}
\hline
client & E[D] & max(D) & min(D) & Var(D) & E[s] & max(s) & Var(s) & plr\\
\hline
500kbs 0.05 & 0.044 & 0.080 & 0.000 & 0.000 & 0.153 & 2.000 & 0.166 & 0.18 \\
\hline
500kbs 0.15 & 0.060 & 0.660 & 0.000 & 0.007 & 0.339 & 6.000 & 1.021 & 0.54 \\
\hline
100kbs 0.05 & 2.050 & 3.070 & 0.520 & 0.530 & 5.085 & 26.000 & 61.286 & 0.70 \\
\hline
100kbs 0.15 & 1.100 & 2.120 & 0.030 & 0.323 & 3.237 & 26.000 & 38.943 & 0.76 \\
\hline
\end{tabular}
\caption{Task 6: Receive statistics for camera sensor data (weakly reliable transport, no congestion control)}
\label{table:task6-camera-b}
\end{table}


\begin{table}
\centering
\begin{tabular}{|c|c|c|c|c|c|c|c|c|}
\hline
client & E[D] & max(D) & min(D) & Var(D) & E[s] & max(s) & Var(s) & plr\\
\hline
500kbs 0.05 & 0.045 & 0.080 & 0.000 & 0.000 & 2.145 & 5.000 & 3.090 & 0.31\\
\hline 
500kbs 0.15 & 0.038 & 0.070 & 0.000 & 0.000 & 0.545 & 6.000 & 1.438 & 0.59\\
\hline
100kbs 0.05 & 2.625 & 3.860 & 0.030 & 1.503 & 5.145 & 35.000 & 91.201 & 0.33\\
\hline
100kbs 0.15 & 2.508 & 4.210 & 0.020 & 2.044 & 4.727 & 34.000 & 88.869 & 0.42\\
\hline
\end{tabular}
\caption{Task 6: Receive statistics for camera sensor data (no reliability guarantees,transport, skipping messages)}
\label{table:task6-camera-c}
\end{table}



\begin{table}
\centering
\begin{tabular}{|c|c|c|c|c|c|c|c|c|}
\hline
client & E[D] & max(D) & min(D) & Var(D) & E[s] & max(s) & Var(s) & plr\\
\hline
500kbs 0.05 & 0.045 & 0.380 & 0.000 & 0.001 & 0.220 & 5.000 & 0.554 & 0.19\\
\hline 
500kbs 0.15 & 0.046 & 0.300 & 0.000 & 0.002 & 0.458 & 8.000 &  2.149 & 0.44\\
\hline
100kbs 0.05 & 1.738 & 3.200 & 0.030 & 0.902 & 5.983 & 28.000 & 78.362 & 0.75\\
\hline
100kbs 0.15 & 1.650 & 2.920 & 0.150 & 0.607 & 5.492 & 28.000 & 74.875 & 0.74\\
\hline
\end{tabular}
\caption{Task 6: Receive statistics for camera sensor data (weakly reliable transport, skipping messages)}
\label{table:task6-camera-d}
\end{table}
